در ادامه بگذارید
$g_k := \nabla E(w^{(k)})$
و
$H_k := \nabla^2 E(w^{(k)})$
و
$p^{(k)} = -H_k^{-1} g_k$
باشد؛ لذا
$w^{(k+1)} = w^{(k)} + p^{(k)}$.

\subsection*{الف}
از تعریف گام نیوتون داریم:
\begin{flalign*}
g_k + H_k p^{(k)} = \nabla E(w^{(k)}) - \nabla^2 E(w^{(k)}) \nabla^2 E(w^{(k)})^{-1}\nabla E(w^{(k)}) = 0. &&
\end{flalign*}
پس:
\begin{flalign*}
\|\nabla E(w^{(k+1)})\|
&= \|\nabla E(w^{(k+1)}) - 0\| \\
&= \|\nabla E(w^{(k+1)}) - (g_k + H_k p^{(k)})\| \\
&= \|\nabla E(w^{(k+1)}) - \nabla E(w^{(k)}) - \nabla^2 E(w^{(k)})p^{(k)}\|. &&
\end{flalign*}

\subsection*{ب}
طبق بسط تیلورِ انتگرالی برای گرادیان (با $x=w^{(k)}$ و $p=p^{(k)}$):
\begin{flalign*}
\nabla E(w^{(k+1)})
&= \nabla E(w^{(k)}+p^{(k)})
= g_k + \int_0^1 \nabla^2 E(w^{(k)}+t p^{(k)})\,p^{(k)}\,dt. &&
\end{flalign*}
بنابراین:
\begin{flalign*}
\nabla E(w^{(k+1)}) - g_k - H_k p^{(k)}
&= \int_0^1 \Big(\nabla^2 E(w^{(k)}+t p^{(k)}) - H_k\Big)p^{(k)}\,dt. &&
\end{flalign*}
اکنون با استفاده از خواص نُرم و انتگرال:
\begin{flalign*}
\|\nabla E(w^{(k+1)})\|
&= \|\nabla E(w^{(k+1)}) - g_k - H_k p^{(k)}\| \\
&\le \int_0^1 \Big\|\big(\nabla^2 E(w^{(k)}+t p^{(k)}) - H_k\big)p^{(k)}\Big\|\,dt \\
&\le \int_0^1 \|\nabla^2 E(w^{(k)}+t p^{(k)}) - H_k\| \cdot \|p^{(k)}\|\,dt. &&
\end{flalign*}

\subsection*{ج}
با فرض $L$-لیپشیتز بودن هسیان:
\begin{flalign*}
\|\nabla^2 E(w^{(k)}+t p^{(k)}) - H_k\|
&\le L\|t p^{(k)}\| = Lt\|p^{(k)}\|. &&
\end{flalign*}
جایگذاری در نامساوی بخش قبل:
\begin{flalign*}
\|\nabla E(w^{(k+1)})\|
&\le \int_0^1 (Lt\|p^{(k)}\|)\|p^{(k)}\|\,dt
= L\|p^{(k)}\|^2\int_0^1 t\,dt
= \frac{L}{2}\|p^{(k)}\|^2. &&
\end{flalign*}
از طرفی
$\|p^{(k)}\|=\|H_k^{-1}g_k\|\le \|H_k^{-1}\|\cdot \|g_k\|$،
پس:
\begin{flalign*}
\|\nabla E(w^{(k+1)})\|
\le \frac{1}{2}L\|H_k^{-1}\|^2\cdot \|\nabla E(w^{(k)})\|^2. &&
\end{flalign*}

\subsection*{د}
اگر $\|H_k^{-1}\|\le 2\|(\nabla^2 E(w^*))^{-1}\|$ در توپِ شعاع $r$ برقرار باشد، آنگاه از نتیجه بخش (ج):
\begin{flalign*}
\|\nabla E(w^{(k+1)})\|
&\le \frac{1}{2}L\big(2\|(\nabla^2 E(w^*))^{-1}\|\big)^2 \|\nabla E(w^{(k)})\|^2 \\
&= 2L\|(\nabla^2 E(w^*))^{-1}\|^2\cdot \|\nabla E(w^{(k)})\|^2. &&
\end{flalign*}
این رابطه نشان می‌دهد اگر به اندازه کافی نزدیک $w^*$ شروع کنیم، نُرم گرادیان‌ها به صفر میل می‌کند (همگرایی به نقطه مانا $w^*$)، و همچنین نرخ همگرایی برای نُرم گرادیان‌ها درجه دو است؛ یعنی در همسایگی $w^*$ داریم:
\begin{flalign*}
\|\nabla E(w^{(k+1)})\| \le C\,\|\nabla E(w^{(k)})\|^2,
\qquad
C=2L\|(\nabla^2 E(w^*))^{-1}\|^2,
&&
\end{flalign*}
پس خطا (به معنای $\|\nabla E(\cdot)\|$) در هر گام تقریباً به توان دو گام قبل کاهش می‌یابد.
